% Proposed replacement for the manuscript's current Section 2.
% This file is a standalone section for integration into report/wcf_main.tex.
% It deliberately leaves K, M, and d generic until the implementation paragraph.

\section{The Wasserstein Causal Forest Model}\label{sec:model}

\subsection{Observed data, representation, and estimands}\label{sec:setup}

Let $O_i=(X_i,A_i,Q_i)$ denote one observational unit, where $X_i\in
\mathbb{R}^{d}$ is a vector of pre-treatment covariates, $A_i\in\{0,1\}$ is
the treatment indicator, and $Q_i=q_K(Y_i)$ is the outcome represented on a
fixed grid $0<u_1<\cdots<u_K<1$. The grid representation is
$q_K(Y)=(Q_Y(u_1),\ldots,Q_Y(u_K))\in\mathcal{Q}_K$, where
$\mathcal{Q}_K=\{q\in\mathbb{R}^K:q_1\leq\cdots\leq q_K\}$. Let
$w_k>0$ and $\sum_k w_k=1$. We use the weighted finite-grid geometry
\[
 d_{W,K}(q,q')=\left\{\sum_{k=1}^{K}w_k(q_k-q'_k)^2\right\}^{1/2}.
\]
This is the declared quadrature approximation to one-dimensional $W_2$
\citep{peyre2019optimaltransport}; all
targets in this paper carry the finite-grid label $K$.

Write $Y^a$ for the potential outcome under treatment level $a$ and
$e(x)=\Pr(A=1\mid X=x)$ for the propensity score. The identification
conditions are consistency, conditional exchangeability
$Y^a\perp A\mid X$, and positivity $0<e(X)<1$ almost surely. Under these
conditions the arm-specific conditional law
\[
 P_a^K(x)=\operatorname{Law}\{q_K(Y^a)\mid X=x\}
 =\operatorname{Law}\{q_K(Y)\mid A=a,X=x\}
\]
is identified. The method estimates this law directly and subsequently
integrates measurable grid functionals against it.

For a measurable scalar functional $h:\mathcal{Q}_K\to\mathbb{R}$, define
\[
 \mu_{a,h}(x)=\mathbb{E}\{h(q_K(Y^a))\mid X=x\},\qquad
 \theta_h=\mathbb{E}\{\mu_{1,h}(X)-\mu_{0,h}(X)\},
\]
and, for a pre-treatment moderator $V=g(X)$,
\[
 \tau_h(v)=\mathbb{E}\{\mu_{1,h}(X)-\mu_{0,h}(X)\mid V=v\}.
\]
We refer to $\theta_h$ as a TATE and to $\tau_h(v)$ as a TCATE. The grid
mean, grid standard deviation, grid skewness, and upper-tail mean are obtained
by taking $h$ to be the corresponding functional of $q_K(Y)$. The identity
functional $h(q)=q$ gives the mean-quantile contrast
\[
 \tau_q(x)=\mathbb{E}\{q_K(Y^1)-q_K(Y^0)\mid X=x\},
\]
which is a vector-valued barycenter contrast and is distinct from nonlinear
outcome-level functionals of the barycenter.

The reference-distribution estimands use the scalar functional
\[
 h_{\mathrm{ref}}(q)=d_{W,K}(q,q_\star),
\]
where $q_\star$ is a fixed benchmark quantile vector. Thus
\[
 \theta_{\mathrm{ref}}=\mathbb{E}\{h_{\mathrm{ref}}(q_K(Y^1))
 -h_{\mathrm{ref}}(q_K(Y^0))\},\qquad
 \tau_{\mathrm{ref}}(v)=\mathbb{E}\{h_{\mathrm{ref}}(q_K(Y^1))
 -h_{\mathrm{ref}}(q_K(Y^0))\mid V=v\}.
\]
These are the reference-distance TATE and TCATE, respectively. A negative
value means that treatment moves the outcome law closer to the benchmark on
average within the relevant population. They are differences of expected
distances, not the distance between an expected quantile vector and the
benchmark.

The term endogenous adoption is used here in the observed-confounding sense:
the propensity varies with pre-treatment covariates that also predict the
potential outcome laws. A propensity different from one half is not, by
itself, evidence of endogeneity. For example, a constant propensity of 0.7
would describe unequal but otherwise randomized allocation. Endogenous
adoption arises when $e(X)$ is related to outcome-relevant prognostic
features, such as $\mu_{a,h}(X)$, or when treatment depends on unobserved
potential-outcome information so that conditional exchangeability fails. The
simulation designs use the first case: the adoption index is a function of
covariates that also enter the outcome surfaces. The proposed calibration
therefore targets confounding adjustment under a nonconstant propensity; it
does not identify effects when unmeasured confounding remains after
conditioning on $X$.

\subsection{Particle-law estimation}\label{sec:boosting}

For each arm, the fitted conditional law is an empirical particle law
\[
 \widehat P_{a,M}^K(x)=\frac{1}{M}\sum_{m=1}^{M}\delta_{p_{am}(x)},
 \qquad p_{am}(x)\in\mathcal{Q}_K.
\]
The particles are fitted by gradient boosting \citep{friedman2001greedy} on the collision-smoothed
energy score
\[
 S_\varepsilon(P,q)=\frac{1}{M}\sum_{m=1}^{M}d_{\varepsilon}(p_m,q)
 -\frac{1}{2M^2}\sum_{m,l=1}^{M}d_{\varepsilon}(p_m,p_l),
 \qquad
 d_{\varepsilon}(a,b)=\{d_{W,K}(a,b)^2+\varepsilon^2\}^{1/2}-\varepsilon.
\]
The first term attracts the particle cloud to the observed quantile vector and
the second term preserves dispersion. Each update fits a multi-output tree to
the negative energy gradient, accepts a backtracked step only when the
training score decreases, and projects the updated particles onto
$\mathcal{Q}_K$ by weighted isotonic regression \citep{barlow1972isotonic}. The particle representation
estimates an arm-specific conditional law. It does not identify a coupling of
$Y^0$ and $Y^1$, and differences between matched particles from the two arms
are not individual treatment effects.

The final learner uses a shared covariate partition. Candidate splits are
scored with the pooled multi-output squared-error criterion and are retained
only when each child contains the required number of observations from each
arm. For the contrast-regularized update used in WCF-DR, let $g_a$ be the
mean gradient in a leaf for arm $a$, let $n_a$ be its arm count, let
$\bar g=(n_0g_0+n_1g_1)/(n_0+n_1)$, and let $\delta=g_1-g_0$. With
$n_{\mathrm{eff}}=n_0n_1/(n_0+n_1)$ and
$c_\lambda=n_{\mathrm{eff}}/(n_{\mathrm{eff}}+\lambda)$, the leaf updates are
\[
 \widehat g_0=\bar g-\frac{n_1}{n_0+n_1}c_\lambda\delta,
 \qquad
 \widehat g_1=\bar g+\frac{n_0}{n_0+n_1}c_\lambda\delta.
\]
This shrinks the arm gap while preserving the count-weighted pooled gradient.
The strength $\lambda$ is selected from a fixed candidate set by held-out
energy risk computed only on each validation unit's observed arm. The method
uses a common pooled initialization, so the initial contrast does not simply
reproduce the difference in the two arms' empirical covariate distributions.

The notation $d$, $K$, and $M$ is generic. The simulations use six covariates
in the income track, $K=25$ grid levels, and $M=10$ particles. These are
implementation choices rather than restrictions of the method. Their effects
are assessed in the sensitivity study described in the accompanying
experimental plan.

\subsection{Doubly robust calibration of scalar functionals}\label{sec:dr}

The particle law and the functional calibration layer target different
objects. For a declared functional $h$, let $\widehat m_{a,h}^{(-k)}(X_i)$
be the out-of-fold mean of $h$ over the fitted arm-$a$ particles, and let
$\widetilde e_i$ be a cross-fitted propensity estimate clipped to a declared
overlap interval. With $H_i=h(q_K(Y_i))$, define
\[
 \widehat\phi_i(h)=\widehat m_{1,h}^{(-k)}(X_i)
 -\widehat m_{0,h}^{(-k)}(X_i)
 +\frac{A_i}{\widetilde e_i}
   \{H_i-\widehat m_{1,h}^{(-k)}(X_i)\}
 -\frac{1-A_i}{1-\widetilde e_i}
   \{H_i-\widehat m_{0,h}^{(-k)}(X_i)\}.
\]
The marginal estimator is the sample mean of $\widehat\phi_i(h)$ and the
moderator-specific estimator is its sample mean within the observed bin
$V=v$. In the implementation, WCF-DR applies this layer to the four grid
functionals and $h_{\mathrm{ref}}$. It replaces their reported contrast
columns, while the particle laws and all law-level scores remain those of the
particle booster.

For fixed nuisance limits $m_a^\dagger$ and $e^\dagger$, the conditional bias
of the corresponding population score is
\[
 (e^\dagger-e(X))\left[
 \frac{m_1^\dagger(X)-\mu_{1,h}(X)}{e^\dagger}
 +\frac{m_0^\dagger(X)-\mu_{0,h}(X)}{1-e^\dagger}
 \right].
\]
Consequently, the score has the standard doubly robust property at the level
of the scalar functional: the bias vanishes when the propensity is correct or
when both arm-specific conditional means are correct, subject to positivity,
integrability, and the usual conditions controlling nuisance estimation and
cross-fitting. This property does not make the fitted conditional law doubly
robust, and it does not by itself provide an asymptotic theorem for the
finite-particle learner.

For the simulation study, $\widetilde e$ is obtained from a cross-fitted
logistic regression and clipped to $[0.02,0.98]$. This choice is a simple
working model, not a requirement of the estimator. A calibrated binary
classification method, such as a suitably cross-fitted forest or boosted
classifier, can be substituted when the assignment mechanism is nonlinear.
The overlap interval must then be reported, and clipping can invalidate the
propensity-correctness branch if the true propensity lies outside the clipped
interval. In particular, the IC3 data-generating mechanism allows propensities
down to 0.01 and up to 0.99, so the fixed $[0.02,0.98]$ clip should not be
described as exactly correctly specified in those tails. The outcome-model
branch remains the relevant robustness route there if its conditional means
are consistently estimated.

\subsection{Optional two-part law for a degenerate component}\label{sec:zipt}

Some distribution-valued outcomes have a structural component in which the
entire unit-level outcome law is degenerate at a point. Examples include a
regional expenditure or benefit distribution that is identically zero when
there are no participating units, or a claim-count distribution for a unit
with no claims. This is different from a continuous distribution that merely
has a lower tail near zero. For such data, the optional two-part extension
estimates
\[
 \pi_a(x)=\Pr\{q_K(Y^a)\neq\mathbf{0}\mid X=x\}
\]
with an arm-specific binary classifier, fits the particle booster on the
nondegenerate rows, and assembles
\[
 \widehat P_a^{\mathrm{2p}}(x)=
 \{1-\widehat\pi_a(x)\}\delta_{\mathbf{0}}
 +\widehat\pi_a(x)\widehat P_{a,M}^{+,K}(x).
\]
The extension gives the degenerate component its own estimable probability
and lets every downstream law functional integrate over both parts. The plain
particle booster can place a finite amount of mass near zero, but it has no
separate component-probability parameter and its particle masses are tied to
$1/M$; the two-part assembly makes that structural mass explicit. In this
paper the extension is an optional model component for zero-inflated outcome
laws, not an assumption about the income-track data-generating process.
